\documentclass{scrarticle}

% --------------------------------------------------------------------------------
% ---- Basic math packages

\usepackage{mathpartir}
\usepackage{amsmath,mathtools} % load mathtools before unicode math

% --------------------------------------------------------------------------------
% ---- Fonts

\usepackage{fontspec}
\setmainfont{Libertinus Serif}
\setsansfont{Libertinus Sans}
% \setmonofont{Libertinus Mono} % really ugly
% For math fonts one seems to also need unicode-math
\usepackage{unicode-math}
\setmathfont{Latin modern math}
\setmathfont{GFSNeohellenicMath}[range={\boxbslash}]
\setmathfont{XITS Math}[range={\smalltriangleright}]

% --------------------------------------------------------------------------------
% ---- Bibliography

\usepackage[backend=biber]{biblatex}
\addbibresource{references.bib}

% --------------------------------------------------------------------------------
% ---- Technical packages

\usepackage{xcolor}
\definecolor{darkred}{rgb}{0.45,0,0}
\definecolor{darkgreen}{rgb}{0,0.45,0}
\usepackage[pdfusetitle,colorlinks,citecolor=darkgreen,menucolor=darkred]{hyperref}
\usepackage{xparse}
\usepackage{xspace}

% --------------------------------------------------------------------------------
% ---- Lists

\usepackage[inline]{enumitem}
\setlist[enumerate,1]{label={(\arabic*)}}

% --------------------------------------------------------------------------------
% ---- Environments

\usepackage{keytheorems}
\usepackage{zref-clever}

\renewkeytheoremstyle{plain}{bodyfont=\normalfont}

\newkeytheorem{axiom}[style=plain]
\newkeytheorem{theorem}[parent=section,style=plain]
\newkeytheorem{conjecture}[parent=section,style=plain]
\newkeytheorem{lemma}[sibling=theorem,style=plain]
\newkeytheorem{proposition}[sibling=theorem,style=plain]
\newkeytheorem{corollary}[sibling=theorem,style=plain]
\newkeytheorem{definition}[sibling=theorem,style=plain,qed=\(\varlrtriangle\)]
\newkeytheorem{extension}[sibling=theorem,style=plain]
\newkeytheorem{construction}[sibling=theorem,style=plain]
\newkeytheorem{notation}[sibling=theorem,style=plain]
\newkeytheorem{example}[sibling=theorem,style=plain,qed=\(\varlrtriangle\)]
\newkeytheorem{observation}[sibling=theorem,style=plain,name=Observation,refname={observation,observations},Refname={Observation,Observations},qed=\(\varlrtriangle\)]
\newkeytheorem{remark}[sibling=theorem,style=remark,qed=\(\varlrtriangle\)]

\numberwithin{equation}{theorem}


% --------------------------------------------------------------------------------
% ---- Math macros

\ExplSyntaxOn
\DeclareExpandableDocumentCommand{\IfNoValueOrEmptyTF}{mmm}
{
 \IfNoValueTF{#1}
  {#2} % true
  {\tl_if_empty:nTF {#1} {#2} {#3}} % false
}
\ExplSyntaxOff

% ---- Basics
\NewDocumentCommand\co{}{\mathchoice{\colon}{\colon}{\mkern-1mu\colon\mkern-5mu}{\colon}}
\DeclarePairedDelimiter\paren{(}{)}
\DeclarePairedDelimiter\card{\lvert}{\rvert}
\NewDocumentCommand\N{s}{\IfBooleanTF{#1}{\mathsf{N}}{\mathbb{N}}}

\DeclarePairedDelimiterX\setImpl[1]{\{}{\}}{#1}
\NewDocumentCommand\set{sO{}>{\TrimSpaces}m>{\TrimSpaces}o}{
  \IfNoValueOrEmptyTF{#3}{\emptyset}{
    \IfValueTF{#4}{
      \setImpl[#2]{#3 \mathbin{\delimsize|} \IfBooleanTF{#1}{\text{#4}}{#4}}
    }{
      \setImpl[#2]{#3}
    }
  }
}

% ---- Category theory
\DeclareMathOperator\id{id}
\NewDocumentCommand\category{m}{\mathbf{#1}}
\NewDocumentCommand\cat{m}{\mathscr{#1}}
\NewDocumentCommand\PSh{om}{\category{PSh}\IfValueT{#1}{_{#1}}(#2)}
\NewDocumentCommand\Cat{}{\category{Cat}}
\NewDocumentCommand\Set{}{\category{Set}}
\NewDocumentCommand\CwF{}{\category{CwF}}
\NewDocumentCommand\yo{m}{\mathbf{y}#1} % よ
\NewDocumentCommand\Nerve{m}{N_{#1}}
\NewDocumentCommand\Lan{mo}{{#1}_!}
\NewDocumentCommand\Res{m}{{#1}^*}
\NewDocumentCommand\Ran{mo}{{#1}_*}
\NewDocumentCommand\op{}{\mathrm{op}}
\NewDocumentCommand\comma{}{\mathbin{\downarrow}}
\NewDocumentCommand\elems{om}{\smallint\IfValueTF{#1}{\nolimits\!\!_{#1}}{\!\!}#2}
\NewDocumentCommand\Arr{m}{{#1}^{\to}}
\NewDocumentCommand\RightCone{m}{{#1}^{\smalltriangleright}}
\NewDocumentCommand\fibSlice{}{\mathbin{/\!\!/}}
\NewDocumentCommand\cc{}{\mathbin{\circ}}
\DeclarePairedDelimiter\pair{\langle}{\rangle}
\DeclarePairedDelimiter\copair{[}{]}
\DeclareMathOperator*\Pullback{\times}
\DeclareMathOperator*\Pushout{+}

\NewDocumentCommand\marked{}{\mathsf{m}}
\NewDocumentCommand\separated{}{\mathsf{sep}}
\NewDocumentCommand\Spine{}{\mathsf{spine}}
\NewDocumentCommand\inner{}{\mathsf{inner}}

\NewDocumentCommand\simplex{m}{\Delta^{#1}}
\NewDocumentCommand\spine{m}{\mathrm{Sp}^{#1}}
\NewDocumentCommand\horn{mm}{\Lambda^{#1}_{#2}}
\NewDocumentCommand\SemiSimplex{s}{\IfBooleanTF{#1}{\Delta_{+,\marked}}{\Delta_+}}
\NewDocumentCommand\semiSimplex{m}{\Delta_{+}^{#1}}


\NewDocumentCommand\fib{}{\twoheadrightarrow}
\NewDocumentCommand\cof{}{\rightarrowtail}
\NewDocumentCommand\ff{}{\hookrightarrow}

\NewDocumentCommand\Leibniz{m}{\mathbin{\widehat{#1}}}
\NewDocumentCommand\Match{mm}{\mathrm{M}_{#1}{#2}}
\NewDocumentCommand\sk{m}{\mathrm{sk}_{#1}}
\NewDocumentCommand\cosk{m}{\mathrm{cosk}_{#1}}
\NewDocumentCommand\res{m}{\mathrm{res}_{#1}}
\DeclareMathOperator*\colim{colim}


% ---- Type theory
\NewDocumentCommand\Univ{smo}{\IfBooleanTF{#1}{\dot{\mathcal{#2}}}{\mathcal{#2}}\IfValueTF{#3}{_{#3}}{}} % TODO: remove old U, V, W macros
\NewDocumentCommand\HS{}{\mathrm{HS}}
\NewDocumentCommand\Rd{}{\mathrm{Rd}}
\NewDocumentCommand\Lv{}{\mathrm{Lv}}
\NewDocumentCommand\Fib{}{\mathrm{Fib}}

\NewDocumentCommand\U{s}{\IfBooleanTF{#1}{\mathcal{U}_{\mathsf{i}}}{\mathcal{U}}}
\NewDocumentCommand\V{}{\mathcal{V}}
\NewDocumentCommand\W{}{\mathcal{W}}
\DeclarePairedDelimiter\trunc{\lVert}{\rVert}

\NewDocumentCommand\Ty{}{\mathrm{Ty}}
\NewDocumentCommand\Prop{}{\mathrm{Prop}}
\NewDocumentCommand\Tm{}{\mathrm{Tm}}
\NewDocumentCommand\p{}{\mathsf{p}}
\NewDocumentCommand\q{}{\mathsf{q}}


% ---- Type theories
\NewDocumentCommand\ETT{}{{\mathbb{E}\textsc{tt}}}

\NewDocumentCommand\UIP{}{\mathsf{UIP}}
\NewDocumentCommand\FE{}{\mathsf{FE}}
\NewDocumentCommand\AC{}{\mathsf{AC}}
\NewDocumentCommand\PT{}{\mathsf{PT}}
\NewDocumentCommand\GC{}{\mathsf{GC}}
\NewDocumentCommand\EM{}{\mathsf{EM}}


\NewDocumentCommand\Id{}{\mathsf{Id}}
\NewDocumentCommand\refl{o}{\mathsf{refl}\IfValueT{#1}{_{#1}}}
\NewDocumentCommand\elim{}{\mathsf{elim}}

% ---- Abbreviations

\ExplSyntaxOn
\NewDocumentCommand\DeclareAbbreviation{mm}{
  \newcommand{#1}{
    \peek_charcode:NTF.{#2}{#2.\@\xspace}
  }
}
\ExplSyntaxOff

\DeclareAbbreviation{\ie}{i.e}
\DeclareAbbreviation{\eg}{e.g}
\DeclareAbbreviation{\cf}{cf}
\DeclareAbbreviation{\etc}{etc}
\DeclareAbbreviation{\resp}{resp}
\DeclareAbbreviation{\etal}{et al}
\DeclareAbbreviation{\wrt}{w.r.t}


% --------------------------------------------------------------------------------
% ---- TikZ

\usepackage{tikz-cd}
% \tikzcdset{arrow style=tikz, diagrams={>={Latex}}}
\tikzcdset{every arrow/.append style = -{>[scale=0.9]}}
\tikzset{
  phantom/.style = {/tikz/commutative diagrams/phantom},
  near/.style = {phantom, very near start},
  pullback/.style = {near, "\lrcorner"},
  pullback-mirrored/.style = {near, "\llcorner"},
  pushout/.style = {near, "\ulcorner"},
  pushout-mirrored/.style = {near, "\urcorner"},
  % fully faithful functors
  ff/.style = {commutative diagrams/hook},
  % Epis and monos
  epi/.style = {commutative diagrams/two heads},
  mono/.style = {commutative diagrams/tail},
  % Model categories
  cof/.style = {{Triangle[open,reversed,scale=0.9]}->},
  fib/.style = {-{Triangle[open,scale=0.9]}},
  weq/.style = {"\sim"'{sloped,font=\tiny,#1}},
  weq'/.style = {"\sim"{sloped,font=\tiny,#1}},
  tcof/.style = {cof,weq={#1}},
  tfib/.style = {fib,weq={#1}},
  % Models of type theory
  % equiv/.style = {weq},
  % disp/.style = {fib}
}
\NewDocumentCommand\mono{}{\rightarrowtail}


% --------------------------------------------------------------------------------
% ---- Main document

\usepackage{todonotes}

\title{Choice without excluded middle?}
\author{Jonas Höfer}
\date{}

\begin{document}
\maketitle

\section{Preliminaries}

We define an extension of extensional type theory which we will still just denote \(\ETT\).
The sorts of the theory are given by a sort of type and a representable sort of terms
\begin{mathpar}
  \vdash\Ty\co\mdlgwhtsquare,
  \and
  A\co\Ty\vdash\Tm(A):*.
\end{mathpar}
We include basic type formers, such as natural numbers and empty, unit, coproduct, dependent function, and dependent pair types.
These are formalized the usual way\footnote{We make use of the usual shorthands in the logical framework. For example, isomorphisms unfold to two operations and two equalities.}, \eg,
\begin{mathpar}
  A \co \Ty, B \co \Tm(A) \to \Ty \vdash \Pi_{A}B \co \Ty,
  \and
  A \co \Ty, B \co \Tm(A) \to \Ty \vdash \paren[\big]{ \paren{u \co \Tm(A) } \to \Tm(B(u)) } \cong \Tm(\Pi_A B).
\end{mathpar}
For readability, we will omit \(\Tm\) when declaring more rules for this theory.
Since we consider extensional type theory, we include the reflection rule for equality
\begin{mathpar}
  A\co\Ty, u, v \co A\vdash \Id_A(u,v) \co \Ty,
  \and
  A\co\Ty, u, v\co A, p \co \Id_A(u,v) \vdash u = v,
  \and
  A\co\Ty, u\co A, p \co \Id_A(u,u) \vdash p = \refl[u],
  \and
  A\co\Ty, u\co A \vdash \refl[u] \co \Id_A(u,u).
\end{mathpar}
As a consequence, we have that \(\ETT \vdash \UIP\) and \(\ETT \vdash \FE\).
To state forms of the axiom of choice, we include a form of truncation.
We axiomatize them as a reflection into subsingleton types.
We write \(\Prop\) as a shorthand for those \(A \co \Ty\) with \(u, v \co A \vdash u = v\).
\begin{mathpar}
  A\co\Ty \vdash \trunc{A} \co \Prop,
  \and
  A\co\Ty,u,v\co\trunc{A} \vdash u = v,
  \and
  A\co\Ty, u\co A \vdash \eta_A(u) \co\trunc{A},
  \and
  A\co\Ty, P\co\Prop^{\trunc{A}}, u\co\paren[\big]{v \co A } \to P(\eta_A(v)),w\co\trunc{A}\vdash \elim(u,w) \co P(w),
  \\
  A\co\Ty, P\co\Prop^{\trunc{A}}, u\co\paren[\big]{v \co A } \to P(\eta_A(v)),w\co A\vdash \elim(u, \eta_A(w)) = u(w).
\end{mathpar}
Using the truncation we can now state different versions of the axiom of choice.

\begin{definition}
  Over \(\ETT\), we define the following principles:
  \begin{description}[leftmargin=0cm]
    \item[\(\AC\)] The \emph{axiom of choice} states that every total relation contains a function, \ie, for all \(A, B \co \Ty\) and \(a \co A, b \co B \vdash R(a, b) \co \Prop\) there is some \(u \co \prod_{a \co A} \trunc{\Sigma_{B} R(a)} \vdash \kappa_A(u) \co \trunc{\prod_{a \co A}\sum_{B} R(a)}\).
    \item[\(\PT\)] The principle of \emph{projective types} states that every type is projective, \ie, for all \(A \co \Ty\) and \(a \co A \vdash P(a) \co \Ty\) there is some \(u \co \Pi_{a\co A}\trunc{P(a)} \vdash \kappa_A(u) \co \trunc{\Pi_AP}\).
    \item[\(\GC\)] The principle of \emph{global choice} states that the support of every type splits, \ie, for all \(A \co \Ty\) there is an element \(u \co \trunc{A} \vdash \varepsilon_A(u) \co A\). \qedhere
  \end{description}
\end{definition}

Over \(\ETT\) the principles \(\AC\) and \(\PT\) are equivalent.
The principle \(\GC\) is in general stronger.
In the presence of a universe we can merely extract a choice function restricted to this universe by using that \(\prod_{A \co \U, p \co \trunc{A}} \trunc{A}\) is inhabited by the identity.

\begin{proposition}
  \(\ETT \vdash \AC \leftrightarrow \PT\) and \(\ETT + \GC \vdash \AC, \PT\).
\end{proposition}
\begin{proof}
  We only show the first claim.
  \todo[inline]{Do; by singleton contraction and UIP}
\end{proof}

Our goal is to relate the above to (restricted versions of) the law of the excluded middle.

\begin{definition}
  Over \(\ETT\), we define the following principles:
  \begin{description}[leftmargin=0cm]
    \item[\(\EM\)] The \emph{law of the excluded middle} states that every proposition is decidable, \ie, for all \(A \co \Prop\) we have \(\vdash \delta_A \co A + \neg A\).
    \item[\(\EM_=\)] The principle of \emph{decidable equality} is \(\EM\) restricted to identity types, \ie, for all \(A \co \Ty\) and \(u, v \co \Ty\) we have \(\vdash \delta_{\Id_A}(u, v) \co \Id_A(u, v) + \neg \Id_A(u, v)\). \qedhere
  \end{description}
\end{definition}

By \(\UIP\) we have that \(\ETT + \EM \vdash \EM_=\).
An observation of Dominik Kirst is that over \(\ETT\) the principle \(\AC\) implies \(\EM_=\) by adapting Diaconescu's theorem.
It does not seem to imply full \(\EM\) without propositional extensionality.
Our goal is to show that this is indeed the case.

\begin{theorem}
  \(\ETT + \AC \vdash \EM_=\).
\end{theorem}
\begin{proof}
  Let \(x, y \co A\) and define the type
  \(
    I \coloneqq \sum_{z \co A} \trunc[\big]{ \Id_A(z, x) + \Id_A(z, y) }
  \)
  and the function \(s \coloneqq \copair{x, y} \co 2 \to I\).
  
  We show that \(s\) is surjective. 
  Given \((z, p) \co I\) we want to show that the fiber over it is merely inhabited.
  Thus, we can assume that we have an element of \(\Id_A(z, x) + \Id_A(z, y)\).
  Without loss of generality, assume we have an element of \(\Id_A(z, x)\).
  Hence, \(z = x\) by induction and an element in the fiber over \((z, p)\) is \(0 \co A\).

  By \(\AC\) the function \(s\) splits.
  Hence, \(I\) is a retract of \(2\) and therefore finite and in particular either of cardinality \(1\) or \(2\).
  If \(I \cong 1\) then \(s(x)\) and \(s(y)\) are equal.
  If \(I \cong 2\) then \(s(x)\) and \(s(y)\) cannot be equal since this would imply that \(s\) is not injective, but it must be bijective.
\end{proof}


\section{Hopefully a counter model}

Let \(\omega\) denote the partially ordered set \(\set{0 < 1 < \ldots}\) viewed as an inverse category.
This means \(\omega^{\op}\) is the category \(\ldots \to 1 \to 0\).
We equip \(\omega^{\op}\) with the Grothendieck topology in which every morphism is covering. 
We will consider separated presheaves for this topology.
A presheaf is separated if there exists at most one patching, \ie, a separated presheaf \(A\) on \(\omega^{\op}\) is given by an ascending chain of monomorphisms \(A_0 \mono A_1 \mono \ldots\).

The category of separated presheaves is a quasitopos~\cite[A2.6]{Johnstone2002Elephant1} which we denote \(\cat{E}\).
It is a full subcategory of the presheaf topos \([\omega, \Set]\).
We recall a standard result about \(\cat{E}\) being reflective in \([\omega, \Set]\).
Furthermore, we note that the unit of the reflector is a levelwise epimorphism which will be crucial later.

\begin{proposition}\label{prop:quasi-topos-is-reflective}
  The inclusion \(U \co \cat{E} \ff [\omega, \Set]\) has a left adjoint \(L \co [\omega, \Set] \to \cat{E}\).
  Furthermore, the unit \(\eta \co \id \to UL\) is valued in levelwise epimorphisms.
\end{proposition}
\begin{proof}
  This is~\cite[page~43,paragraph~3]{Vistoli2008}.
  Concretely, for \(X \in [\omega, \Set]\) the separated presheaf \(LX\) is given at \(n\in\omega\) by the quotient of \(X^n\) that identifies two local sections \(u, v \in X^n\) if there exists some \(m \geq n\) such that \(u|_m = v|_m\).
\end{proof}

\todo[inline]{I think we also need this to be stable}

\begin{example}\label{ex:coequalizer-are-not-levelwise-in-quasi-topos}
  The forgetful functor \(\cat{E} \to [\omega, \Set]\) does not create coequalizer.
  Consider the following example where all involved presheaves are constant on \(\omega_{\geq 1}\).
  \[\begin{tikzcd}
    0 && 2 && 2 \\
    \\
    {2 \times 2} && 2 && 1 \\
    \\
    \ldots && \ldots && \ldots
    \arrow[shift left=2, from=1-1, to=1-3]
    \arrow[shift right=2, from=1-1, to=1-3]
    \arrow[from=1-1, to=3-1]
    \arrow[dashed, two heads, from=1-3, to=1-5]
    \arrow[equals, from=1-3, to=3-3]
    \arrow[two heads, from=1-5, to=3-5]
    \arrow[shift left=2, from=3-1, to=3-3]
    \arrow[shift right=2, from=3-1, to=3-3]
    \arrow[equals, from=3-1, to=5-1]
    \arrow[dashed, two heads, from=3-3, to=3-5]
    \arrow[equals, from=3-3, to=5-3]
    \arrow[equals, from=3-5, to=5-5]
  \end{tikzcd}\]
  The top row is trivially a coequalizer.
  The second and all following rows are exactly the image factorization.
\end{example}

\begin{proposition}\label{prop:forgetful-functor-creates}
  The forgetful functor \(U \co \cat{E} \ff [\omega, \Set]\) creates
  \begin{enumerate*}
    \item\label{prop:forgetful-functor-creates:limits} limits,
    \item\label{prop:forgetful-functor-creates:coproducts} coproducts,
    \item\label{prop:forgetful-functor-creates:epimorphisms} epimorphisms,
    \item\label{prop:forgetful-functor-creates:images} image factorizations,
    \item\label{prop:forgetful-functor-creates:exponentials} exponentials, and
    \item\label{prop:forgetful-functor-creates:pushforwards} pushforwards.
  \end{enumerate*} 
\end{proposition}
\begin{proof}
  \begin{description}[leftmargin=0cm]
    \item[\ref{prop:forgetful-functor-creates:limits}] By \zcref{prop:quasi-topos-is-reflective}, the inclusion is monadic and therefore creates limits.
    \item[\ref{prop:forgetful-functor-creates:coproducts}] Coproducts of monomorphisms in \(\Set\) are monomorphisms.
    \item[\ref{prop:forgetful-functor-creates:epimorphisms}] Epimorphisms in \([\omega,\Set]\) are levelwise.
      Since two morphisms in \(\cat{E}\) agree if they are equal levelwise, we have that levelwise epimorphisms are also epi in \(\cat{E}\).
    \item[\ref{prop:forgetful-functor-creates:images}] Consider the factorization of a morphism \(f \co A \to B\) between objects of \(\cat{E}\) in \([\omega, \Set]\):
      \[\begin{tikzcd}
        {A_0} && {\mathrm{Img}(f_0)} && {B_0} \\
        \\
        {A_1} && {\mathrm{Img}(f_1)} && {B_1\rlap{.}}
        \arrow[two heads, from=1-1, to=1-3]
        \arrow[tail, from=1-1, to=3-1]
        \arrow[tail, from=1-3, to=1-5]
        \arrow[dashed, from=1-3, to=3-3]
        \arrow[tail, from=1-5, to=3-5]
        \arrow[two heads, from=3-1, to=3-3]
        \arrow[tail, from=3-3, to=3-5]
      \end{tikzcd}\]
      By the left-cancellation property of monomorphisms, the dashed map is a also mono.
    \item[\ref{prop:forgetful-functor-creates:exponentials}]
    \item[\ref{prop:forgetful-functor-creates:pushforwards}]
  \end{description}
\end{proof}

\begin{example}\label{ex:epimorphisms-are-not-levelwise-in-quasi-topos}
  Not every epimorphism in \(\cat{E}\) is levelwise.
  Consider the morphism \(i\) in \([2, \Set]\) between monomorphisms below.
  The example can be extended to a one in \(\cat{E}\), as in \zcref{ex:coequalizer-are-not-levelwise-in-quasi-topos}.
  \[\begin{tikzcd}
    0 && 1 && {A_0} \\
    \\
    2 && 1 && {A_1}
    \arrow["{i_0}", from=1-1, to=1-3]
    \arrow[tail, from=1-1, to=3-1]
    \arrow["{f_0}", shift left=2, from=1-3, to=1-5]
    \arrow["{g_0}"', shift right=2, from=1-3, to=1-5]
    \arrow[tail, from=1-3, to=3-3]
    \arrow["{a}", tail, from=1-5, to=3-5]
    \arrow["{i_1}", two heads, from=3-1, to=3-3]
    \arrow["{f_1}", shift left=2, from=3-3, to=3-5]
    \arrow["{g_1}"', shift right=2, from=3-3, to=3-5]
  \end{tikzcd}\]
  Suppose \(f \cc i = g \cc i\).
  Since \(i_1\) is epi, we have \(f_1 = g_1\).
  Thus, \(a \cc f_0 = a \cc g_0\).
  Since \(a\) is mono, we have \(f_0 = g_0\).
  Consequently, the forgetful functor does not preserve epimorphisms.
\end{example}


\subsection{An interpretation of type theory}

Fix a Grothendieck universe \(\V\).
We write \(\V_{\HS}\) for the Hofmann-Streicher lifting of the universe and \(\V_{\HS,\separated} \subseteq \V_{\HS}\) for the subuniverse on separated presheaves.
These universes are not separated since the functorial actions on \(n < m\) is given by restriction along \(\omega_{\geq m} \ff \omega_{\geq n}\).

\begin{observation}
  The category \(\omega^{\op}\) has a terminal object.
  Hence, \(\Gamma \co \PSh{\omega^{\op}} \to \Set\) is given by evaluation.
  Consequently, in the sequence of adjoints \(\pi_0 \dashv \Delta \dashv \Gamma \dashv \nabla\) the functors \(\Delta\) and \(\nabla\) are fully faithful since they are Kan extensions along a fully faithful functor.
  
  Thus, \(\Set\) is embeds into \([\omega, \Set]\) in two ways.
  Furthermore, \(\Delta\) factors over the subcategory \(\cat{E}\) since all restriction operations of a constant presheaf are in particular mono.
  Hence, the adjoint triple \(\pi_0 \dashv \Delta \dashv \Gamma\) descends to one between \(\Set\) and \(\cat{E}\), witnessing \(\Set\) to be reflective and coreflective in \(\cat{E}\).
  We obtain the usual global sections comonad \(\Delta \cc \Gamma\) on \(\cat{E}\). 
\end{observation}

\begin{construction}
  The universes \(\dot{\U} \to \U\) and \(\dot{\U}_{\separated} \to \U_{\separated}\) are defined by pulling back in \([\omega, \Set]\) along the counit \(\varepsilon \co \Delta\Gamma \Rightarrow \id\) of the global sections comonad:
  \[\begin{tikzcd}
    {\dot{\U}_{\separated}} && {\dot{\V}_{\HS, \separated}} & \\
    & {\dot{\U}} && {\dot{\V}_{\HS}} \\
    {\U_{\separated}} && {\V_{\HS, \separated}} \\
    & \U && {\V_{\HS}\rlap{.}}
    \arrow[from=1-1, to=1-3]
    \arrow[from=1-1, to=2-2]
    \arrow[from=1-1, to=3-1]
    \arrow[from=1-3, to=2-4]
    \arrow[from=1-3, to=3-3]
    \arrow[from=2-2, to=2-4]
    \arrow[from=2-2, to=4-2]
    \arrow[from=2-4, to=4-4]
    \arrow["\varepsilon"', from=3-1, to=3-3]
    \arrow[from=3-1, to=4-2]
    \arrow[from=3-3, to=4-4]
    \arrow["\varepsilon"', from=4-2, to=4-4]
  \end{tikzcd}\]
  The left square is a pullback and the universe \(\dot{\U}_{\separated} \to \U_{\separated}\) lays in \(\cat{E}\).
  Furthermore, the above components of the counit are epimorphisms.
\end{construction}
\begin{proof}
  The right face of the cube is a pullback square by definition.
  Hence, by pasting, the left square is a pullback as-well.

  % surj: use that we can just copy the first stage
  % in E: use that the sep. one is fiberwise sep. Hence, pulling back to something in E yields something in E

  \todo[inline]{Finish}
\end{proof}


% We define a universe \(\U\) in \(\cat{E}\) by setting \(\U\) be the constant presheaf on \(\V\)-small objects of \(\cat{E}\) and setting \(\dot{\U}_n \coloneqq \coprod_{A\in \U_n} A_n\) with the evident functorial action. 
% There is an evident projection \(\p \co \dot{\U} \to \U\).
% Note that both \(\U\) and \(\dot{\U}\) lay in \(\cat{E}\).

% We define a model of the theory \(\ETT\).
% The category of contexts is given by \(\cat{E}\).
% The presheaf of types is the one represented by \(\U\), \ie, we set \(\Ty_{\cat{E}} \coloneqq \cat{E}(-, \U)\).
% Given \(A \co \Gamma \to \U\) the context extension \(\p_A \co \Gamma.A \to \Gamma\) is the chosen pullback of \(\p \co \dot{\U} \to \U\) along \(A \co \Gamma \to \U\).
% This is well-defined by \zcref{prop:forgetful-functor-creates}~\ref{prop:forgetful-functor-creates:limits}.
% This defines a category with attributes.
% By the equivalence with CwFs, this induces a presheaf of terms \(\Tm_{\cat{E}}\) which at \(\Gamma \in \cat{E}\) and \(A \co \Gamma \to \U\) is given by the set \(\Tm_{\cat{E}}(\Gamma, A)\) of sections of \(\p_A\).





\printbibliography

\newpage


\section{Question}

The goal is to produce a model of type theory with a non-univalent universe of propositions, decidable equality, negation of excluded middle, and choice.
This might be possible since in the absence of univalence for the universe of propositions Diaconescu's theorem does not seem to go through.
We will work in a classical meta-theory, \ie, assume the axiom of choice in \(\Set\).

\begin{definition}
  Let \(\mathbb{T}\) be the extension of type theory by uniqueness of identity proofs and function extensionality.

  We write \(\exists_{a \co A} P(a)\) as a shorthand for \(\trunc{\sum_{A}P}\).
\end{definition}

\begin{definition}
  We define the following principles:
  \begin{description}
    \item[PT] For all types \(\vdash A\) and families \(a \co A \vdash P(a)\), given \(a \co A \vdash \trunc{P(a)}\) we have \(\vdash \trunc{ \prod_{a \co A} P(a) }\).
    \item[AC] For all types \(\vdash A, B\) and relations \(a \co A, b \co B \vdash R(A, B) \co \Omega\) such that \(a \co A \vdash \exists_{b \co B} R(a, b)\) we have \(\vdash \exists_{f \co A \to B} \prod_{a \co A} R(a, f(a))\). \qedhere
  \end{description}
\end{definition}


\todo[inline]{Are the above equivalent in our setting?}


\section{Candidate?}

\todo[inline]{Does FE hold in this model?}

Let \(\cat{S}\) be the Sirpinski topos, which we view as covariant presheaves on \([1] = \set{0 < 1}\), \ie, as \(\Arr{\Set}\).
We consider the Grothendieck topology with a single non-trivial sieve, which is generated by \(0 < 1\).
Denote the quasi-topos of separated presheaves for this topology by \(\cat{E}\).

Recall that \(\Set\) has an orthogonal factorization system given by surjections and injections and that this factorization system is stable.

\begin{lemma}\label{prop:separation-is-factorization}
  The separation functor \(L \co \cat{S} \to \cat{E}\) is given by the second factor of the epi-mono factorization in \(\Set\).
\end{lemma}
\begin{proof}
  The catgory \(\cat{E}\) of separated presheaves is the subcategory of the arrow category of \(Set\) on monomorphisms.
  The right factor functor for an orthogonal factorization system is exactly the reflection into right maps.
\end{proof}

We first establish some closure properties of \(\cat{E} \subseteq \cat{S}\).

\begin{lemma}
  The forgetful functor \(\cat{E} \to \cat{S}\) creates the following:
  \begin{enumerate}
    \item limits,
    % \item filtered colimits, % can be read of from the orthogonality class generating this factorization system
    \item truncations,
    \item exponentials,
    \item pushforwards.
  \end{enumerate}
\end{lemma}
\begin{proof}
  The first claim follow from the monadic functors creating limits.
  For the second claim we consider the truncation of an object \(A \in \cat{E}\), \ie, the image of the unique map \(A \to 1\).
  In \(\cat{S}\), this image can be computed as the following coequalizer:
  \[\begin{tikzcd}
    {A_0 \times A_0} && {A_0} && {\trunc{A_0}} && 1 \\
    \\
    {A_1 \times A_1} && {A_1} && {\trunc{A_1}} && {1\rlap{.}}
    \arrow["{\pi_0}", shift left=2, from=1-1, to=1-3]
    \arrow["{\pi_1}"', shift right=2, from=1-1, to=1-3]
    \arrow[tail, from=1-1, to=3-1]
    \arrow[two heads, from=1-3, to=1-5]
    \arrow[tail, from=1-3, to=3-3]
    \arrow[tail, from=1-5, to=1-7]
    \arrow[dashed, from=1-5, to=3-5]
    \arrow[equals, from=1-7, to=3-7]
    \arrow["{\pi_0}", shift left=2, from=3-1, to=3-3]
    \arrow["{\pi_1}"', shift right=2, from=3-1, to=3-3]
    \arrow[two heads, from=3-3, to=3-5]
    \arrow[tail, from=3-5, to=3-7]
  \end{tikzcd}\]
  But now the dashed morphisms is mono by a cancelation property of monomorphisms.
  Hence, it is already in the reflective subcategory.

  For the third claim we have to show that the map \((A^B)_0 \to (A^B)_1\) is mono.
  We have that \((A^B)_0 = \Arr{\Set}(A, B)\) and that \((A^B)_1 = A_1^{B_1}\).
  Consider any morphism from \(A\) to \(B\):
  \[\begin{tikzcd}
    {A_0} && {B_0} \\
    \\
    {A_1} && {B_1\rlap{.}}
    \arrow["{f_0}", from=1-1, to=1-3]
    \arrow[tail, from=1-1, to=3-1, "a"]
    \arrow[tail, from=1-3, to=3-3, "b"]
    \arrow["{f_1}"', from=3-1, to=3-3]
  \end{tikzcd}\]
  The component \(f_0\) is uniquely determined since \(f_0 \cc b = f_1 \cc a\) and \(b\) is mono.
  Hence, there exists at most one factorization of \(f_1 \cc a\) over \(b\) making the square commute.

  For the fourth claim we use the characterization of a pushforward as the fiber of an exponential.
\end{proof}

\begin{corollary}
  Monomorphisms and truncations \(\eta_A \co A \to \trunc{A}\) in \(\cat{E}\) are levelwise.
\end{corollary}

\begin{corollary}
  The quasi topos \(\cat{E}\) saitisfies function exponentiality.
\end{corollary}
\begin{proof}
  The statement is computed as in \(\cat{S}\).
\end{proof}

\begin{lemma}
  Coequalizer in \(\cat{E}\) are levelwise epimorphisms.
\end{lemma}
\begin{proof}
  A colimit in a reflective subcategory is computed as the reflection of the colimit in the ambient category.
  By \zcref{prop:separation-is-factorization}, the reflection is given by the epi-mono factorization in \(\Set\).
  Hence, the coequalizer is given by the following monomorphism on the right:
  \[\begin{tikzcd}
    {R_0} && {A_0} && {A_0 / R_0} && {\trunc{A_0 / R_0}} \\
    \\
    {R_1} && {A_1} && {A_1 / R_1} && {A_1 / R_1}
    \arrow[shift left=2, from=1-1, to=1-3]
    \arrow[shift right=2, from=1-1, to=1-3]
    \arrow[tail, from=1-1, to=3-1]
    \arrow[two heads, from=1-3, to=1-5]
    \arrow[tail, from=1-3, to=3-3]
    \arrow["{(\eta_{A/R})_0}", two heads, from=1-5, to=1-7]
    \arrow[from=1-5, to=3-5]
    \arrow[tail, from=1-7, to=3-7]
    \arrow[shift left=2, from=3-1, to=3-3]
    \arrow[shift right=2, from=3-1, to=3-3]
    \arrow[two heads, from=3-3, to=3-5]
    \arrow["{(\eta_{A/R})_1}"', equals, from=3-5, to=3-7]
  \end{tikzcd}\]
  The universal cocone is given by composing the universal cocone with the unit of the monad.
  As the unit is epi (since it is given by the first factor of the factorization and the identity) the claim follows.
\end{proof}

\begin{proposition}
  Levelwise epimorphisms in \(\cat{E}\) split.
\end{proposition}
\begin{proof}
  Consider the following levelwise epimorphism in \(\cat{E}\):
  \[\begin{tikzcd}
    {A_0} && {B_0} \\
    \\
    {A_1} && {B_1\rlap{.}}
    \arrow["{f_0}", two heads, from=1-1, to=1-3]
    \arrow[tail, from=1-1, to=3-1]
    \arrow[tail, from=1-3, to=3-3]
    \arrow["{f_1}", two heads, from=3-1, to=3-3]
  \end{tikzcd}\]
  Since every monomorphism in \(\Set\) is decidable, the right map is equivalent to a coproduct injection.
  By extensivitiy, the bottom map can be written as a coproduct of two maps, both of which are epimorphism.
  Hence, it suffices to consider a square as below on the right:
  \[\begin{tikzcd}
    {B_0} && {A_0} && {B_0} \\
    && {A_0'} \\
    {B_0 + B_1'} && {A_0' + A_1'} && {B_0 + B_1'\rlap{.}}
    \arrow["{s_0}", dashed, from=1-1, to=1-3]
    \arrow["{\iota_0}", tail, from=1-1, to=3-1]
    \arrow["{f_0}", two heads, from=1-3, to=1-5]
    \arrow["i", tail, from=1-3, to=2-3]
    \arrow["{\iota_0}", tail, from=1-5, to=3-5]
    \arrow["{f_1'}"', from=2-3, to=1-5]
    \arrow["{\iota_0}", from=2-3, to=3-3, tail]
    \arrow["{i \cc s_0 + s_1}", dashed, from=3-1, to=3-3]
    \arrow["{f_1' + f_1''}", two heads, from=3-3, to=3-5]
  \end{tikzcd}\]
  Now every splitting \(s_0\) of \(f_0\) induces one of \(f_1'\).
  Thus, by splitting \(f_0\) and \(f_1''\) we obtain a splitting of the bottom map such that the left square commutes by construction.
\end{proof}


\printbibliography%

\end{document}
